\documentclass[11pt]{res}

\usepackage[dvipsnames]{xcolor}

\definecolor{darkblue}{RGB}{46,48,147}
\definecolor{inspirecol}{RGB}{16,23,39}
\definecolor{publonscol}{RGB}{51,102,153}
\definecolor{arxivcol}{RGB}{178,27,27}
\definecolor{orcidlogocol}{HTML}{A6CE39}
\definecolor{webofscicolone}{HTML}{6719FA}
\definecolor{webofscicoltwo}{HTML}{1BC407}
\usepackage{scalerel}
\usepackage{tikz}
\usetikzlibrary{svg.path}

\tikzset{
  orcidlogo/.pic={
    \fill[orcidlogocol] svg{M256,128c0,70.7-57.3,128-128,128C57.3,256,0,198.7,0,128C0,57.3,57.3,0,128,0C198.7,0,256,57.3,256,128z};
    \fill[white] svg{M86.3,186.2H70.9V79.1h15.4v48.4V186.2z}
                 svg{M108.9,79.1h41.6c39.6,0,57,28.3,57,53.6c0,27.5-21.5,53.6-56.8,53.6h-41.8V79.1z M124.3,172.4h24.5c34.9,0,42.9-26.5,42.9-39.7c0-21.5-13.7-39.7-43.7-39.7h-23.7V172.4z}
                 svg{M88.7,56.8c0,5.5-4.5,10.1-10.1,10.1c-5.6,0-10.1-4.6-10.1-10.1c0-5.6,4.5-10.1,10.1-10.1C84.2,46.7,88.7,51.3,88.7,56.8z};
  }
}
\tikzset{
  inspirelogo/.pic={
    \fill[inspirecol] svg{M1224.462,587.075c17.182-120.956,121.118-213.986,246.82-213.986,137.72,0,249.366,111.646,249.366,249.366s-111.646,249.366-249.366,249.366c-125.704,0-229.632-93.030-246.82-213.986h-652.982c0.724,58.434,1.762,114.418,3.116,167.966l491.776,254.098c29.93-41.698,78.71-68.96,133.954-68.96,91.092,0,164.93,73.846,164.93,164.93,0,18.842-3.314,36.872-9.136,53.738l154.382,79.77-32.476,62.852-158.918-82.118c-30.010,31.196-72.090,50.688-118.792,50.688-91.092,0-164.93-73.846-164.93-164.93,0-10.54,1.090-20.824,2.984-30.816l-461.436-238.424c2.326,70.934,5.238,137.26,8.77,198.752l-258.348-34.882c30.134-168.792,45.202-377.622,45.202-626.476,0-221.316-6.026-418.084-18.080-590.306h233.794c-10.050,155.67-15.842,318.332-17.518,487.754l467.602-241.612c-1.894-9.998-2.984-20.276-2.984-30.822,0-91.092,73.846-164.93,164.93-164.93,46.702,0,88.78,19.492,118.792,50.688l158.918-82.118,32.476,62.852-154.382,79.77c5.814,16.866,9.128,34.896,9.128,53.738,0,91.092-73.846,164.93-164.93,164.93-55.244,0-104.016-27.26-133.946-68.952l-496.092,256.336c-0.014,7.578-0.080,15.068-0.080,22.666,0,73.070,0.424,142.848,1.228,209.576v-66.516h653.048zM479.774,1305.593c51.668,0,91.714,16.142,120.13,48.442,28.416,32.292,42.62,70.392,42.62,114.314,0,37.88-12.266,70.816-36.812,98.808-24.54,27.978-57.922,41.984-100.096,41.984-51.668,0-92.152-15.938-121.418-47.792-29.294-31.868-43.922-69.332-43.922-112.376,0-39.614,12.274-73.42,36.812-101.398,24.54-28,58.77-41.984,102.686-41.984z};
  }
}
\tikzset{
  publonslogo/.pic={
    \fill[publonscol] svg{M1582.022,1339.517c-14.724,35.914-35.818,68.878-59.392,99.548-56.188,70.956-130.392,131.16-218.872,155.414-132.966,36.768-280.35,1.894-388.798-81.568-23.304-17.664-43.988-38.348-63.136-60.342-9.172,47.47-12.018,96.022-22.77,143.126-5.508,5.918-10.84,14.774-20.32,13.282-23.582-0.454-45.48-10.386-66.94-19.148-92.452-40.2-184.59-81.028-277-121.176-9.758-4.652-21.958-7.050-28.19-16.764-2.034-13.012-1.894-26.514-0.138-39.57,5.602-12.822,19.558-18.688,30.902-25.6,29.806-16.758,61.922-30.53,88.072-53.072,23.486-19.332,35.906-49.050,38.034-78.95,0.944-9.896,1.266-19.83,1.31-29.806-0.088-362.87,0.052-725.782-0.088-1088.65-2.348-45.846,0.636-92.094-6.73-137.574-3.796-21.408-16.076-45.348-39.476-49.59-35.050-2.706-70.414,1.31-105.414-1.806-9.354-0.358-12.646-10.526-12.062-18.38,0.088-20.868-0.586-41.78,0.358-62.69-0.044-5.734,6.188-8.448,9.398-12.376,65.398,1.712,130.758,6.458,196.154,7.226,125.016,2.436,250.126,0.226,374.922-7.314,9.802-0.044,25.248-0.182,27.414,12.236,1.946,22.316,1.674,44.938,0.182,67.21-0.182,9.896-10.254,14.592-18.834,14.592-27.012,1.398-54.016,0.176-81.028,0.454-15.712,0.41-32.79-1.894-47.104,6.1-14.27,8.668-19.244,26.016-23.084,41.238-11.382,57-7.182,115.492-9.304,173.216-1.404,128.314-3.342,256.636-3.796,384.994,53.116-23.94,110.796-36.404,168.652-41.508,100.498-9.662,204.062,10.21,292.726,58.85,111.696,59.984,197.428,162.334,245.84,278.682,38.348,91.194,55.238,190.734,53.526,289.426-0.33,68-8.462,137.15-35.020,200.294zM1337.674,809.627c-19.916-70.964-60.248-141.056-127.056-176.83-56.364-31.386-124.072-27.546-185.27-15.404-47.514,9.398-90.646,33.828-127.012,65.222-31.114,30.442-42.686,75.564-42.094,117.972,0.592,184.73-0.856,369.554,0.724,554.24,41.684,58.448,109.436,96.248,180.122,105.962,67.116,9.846,138.116-12.244,186.448-60.072,65.492-63.818,97.828-153.292,116.802-240.874,22.674-115.082,29.002-236.17-2.662-350.216z};
  }
}
\tikzset{
  arxivlogo/.pic={
    \fill[arxivcol] svg{M894.142 947.471c0-71.104-26.83-128.746-59.926-128.746s-59.926 57.642-59.926 128.746c0 71.104 26.83 128.746 59.926 128.746s59.926-57.642 59.926-128.746zM1273.782 947.471c0-71.104-26.798-128.746-59.852-128.746s-59.852 57.642-59.852 128.746c0 71.104 26.798 128.746 59.852 128.746s59.852-57.642 59.852-128.746zM1374.216 666.501h-63.43c0-33.36 12.354-63.926 32.432-88.356-21.578-54.14-150.792-121.856-319.218-121.856-168.36 0-297.64 67.716-319.218 121.856 19.998 24.43 32.504 54.996 32.504 88.356h-63.502c0-44.712-39.57-81.072-88.144-81.072v-63.502c24.568 0 47.572 6.218 68.14 16.216 46.578-78.212 190.288-145.356 370.22-145.356 179.998 0 323.642 67.146 370.212 145.356 20.576-9.998 43.572-16.216 68.148-16.216v63.502c-48.574 0-88.144 36.36-88.144 81.072zM1910.36 164.931c-19.002 35.284-35.43 42.358-66.29 51.924-26.5 8.002-73.070 4.286-113.21-11.352-9.142-3.576-17.43-9.362-26.148-14.072-25.994-13.926-41.64-11.212-61.426 9.926-11.36 12.214-47.002 44.434-95.926 90.858 99.854 118.858 160.146 271.858 160.146 439.216 0 167.286-60.284 320.358-160.146 439.216 48.926 46.358 84.494 78.496 95.854 90.778 19.786 21.146 35.43 23.86 61.432 9.86 8.712-4.644 16.998-10.43 26.142-14 40.068-15.646 86.71-19.36 113.21-11.286 30.786 9.494 47.286 16.566 66.29 51.85 11.644 27.216 15.93 51.214 7.432 81.716-8.506 30.362-43.286 49.862-87.004 57.432-35.072 6.144-20.64 11.85-21.284 43.286-0.856 47.14-17.788 84.202-58.002 109.356-30.428 18.856-65.288 20.356-98.998 19.142-46.358-1.5-94.142-24.43-105.86-51.998-13.926-32.57 1.726-53.57 14.504-76.858 16.502-29.784 8.426-59.64-12.932-84.502-23.070-27.070-79.784-85.284-112.070-119.64-114.57 86.858-257.076 139.060-412 139.060-155.004 0-297.502-52.202-412.146-139.14-32.286 34.442-89 92.658-112.070 119.728-21.358 24.862-29.36 54.718-12.932 84.576 12.786 23.216 28.43 44.214 14.57 76.786-11.784 27.64-59.568 50.498-105.926 51.998-33.712 1.288-68.644-0.212-99.072-19.142-40.14-25.074-57.14-62.216-58.002-109.282-0.57-31.502 13.86-37.216-21.284-43.36-43.71-7.504-78.424-27.070-86.996-57.358-8.5-30.5-4.214-54.572 7.432-81.716 18.996-35.284 35.496-42.43 66.282-51.924 26.5-8.074 73.142-4.286 113.284 11.286 9.142 3.642 17.43 9.428 26.142 14 26.002 14 41.574 11.352 61.426-9.786 11.424-12.362 47.002-44.5 95.854-90.858-99.848-118.858-160.132-271.93-160.132-439.216 0-167.358 60.358-320.498 160.146-439.288-48.86-46.35-84.428-78.57-95.854-90.784-19.858-21.138-35.43-23.852-61.426-9.926-8.712 4.718-16.998 10.504-26.142 14.072-40.14 15.638-86.784 19.354-113.284 11.352-30.786-9.568-47.286-16.64-66.282-51.924-11.644-27.216-15.93-51.222-7.432-81.716 8.572-30.354 43.286-49.854 87.004-57.432 35.146-6.136 20.714-11.856 21.284-43.286 0.856-47.14 17.854-84.21 58.002-109.356 30.428-18.856 65.36-20.356 99.072-19.142 46.358 1.426 94.142 24.43 105.926 51.998 13.86 32.504-1.784 53.642-14.57 76.858-16.428 29.718-8.426 59.574 12.932 84.502 23.070 26.998 79.784 85.284 112.070 119.712 114.644-86.93 257.142-139.14 412.146-139.14 154.998 0 297.428 52.21 412 139.14 32.286-34.428 89.074-92.716 112.142-119.712 21.358-24.928 29.36-54.784 12.932-84.502-12.858-23.216-28.43-44.354-14.578-76.858 11.798-27.574 59.574-50.57 105.932-51.998 33.712-1.214 68.644 0.286 99.080 19.142 40.14 25.146 57.066 62.216 57.93 109.356 0.636 31.43-13.788 37.15 21.358 43.286 43.638 7.578 78.424 27.078 86.922 57.432 8.544 30.5 4.184 54.498-7.454 81.716zM1024.074 90.215c-353.572 0-641.214 287.642-641.214 641.214 0 353.5 287.642 641.142 641.214 641.142 353.426 0 641.142-287.642 641.142-641.142 0-353.572-287.714-641.214-641.142-641.214z};
    }
}
\tikzset{
  webofscilogo/.pic={
    \fill[webofscicolone] svg{M166.5 79.5C161.929 87.5 153 107.5 153 113.5C168.5 116.5 200.5 128 229.5 143C252.7 155 287 188.5 295 197.5C301.5 192 312.6 177.7 317 170.5C321.4 163.3 327.5 151 330 139.5C321.5 130.833 304 109.5 263.5 86C236.344 70.2426 203.167 59.8333 185 55C179.5 60 171.071 71.5 166.5 79.5Z};
    \fill[webofscicoltwo] svg{M317.672 225.24C312.909 217.353 300.241 200.421 295 197.5C284.833 209.578 258.755 231.326 231.534 249.355C209.757 263.777 163.256 278.894 151.5 281.5C153.14 289.855 160.766 305.047 164.913 312.396C169.06 319.745 176.834 331.061 185.662 338.844C197.37 335.638 224.525 330.739 264.769 306.804C291.754 290.755 316.999 266.843 330.065 253.328C328.376 246.089 322.434 233.128 317.672 225.24Z};
    \fill[black] svg{M114.56 286.691C123.915 286.651 146.032 284.351 151.307 281.306C146.074 266.151 139.936 232.17 138.399 199.051C137.169 172.556 149.149 125.104 153 113.5C144.864 110.575 126.711 108.352 118.147 108.137C109.582 107.922 95.671 108.801 84.2909 112.439C80.9873 124.314 71.1173 150.534 71.0212 198.081C70.9567 229.963 78.6513 264.426 83.6264 282.857C90.815 285.156 105.205 286.731 114.56 286.691Z};
  }
}

\newcommand\orcidicon[1]{\href{https://orcid.org/#1}{\mbox{\scalerel*{
  \begin{tikzpicture}[yscale=-1,transform shape]
    \pic{orcidlogo};
  \end{tikzpicture}
  }{|}}}}
\newcommand\inspireicon[1]{\href{https://inspirehep.net/authors/#1}{\mbox{\scalerel*{
  \begin{tikzpicture}[yscale=1,transform shape]
    \pic{inspirelogo};
  \end{tikzpicture}
  }{|}}}}
\newcommand\publonsicon[1]{\href{https://publons.com/researcher/#1}{\mbox{\scalerel*{
  \begin{tikzpicture}[yscale=1,transform shape]
    \pic{publonslogo};
  \end{tikzpicture}
  }{|}}}}
\newcommand\arxivicon[1]{\href{https://arxiv.org/a/#1.html}{\mbox{\scalerel*{
  \begin{tikzpicture}[yscale=1,transform shape]
    \pic{arxivlogo};
  \end{tikzpicture}
  }{|}}}}
\newcommand\webofsciicon[1]{\href{https://www.webofscience.com/wos/author/record/#1}{\mbox{\scalerel*{
\begin{tikzpicture}[yscale=1,transform shape]
\pic{webofscilogo};
\end{tikzpicture}
}{|}}}}

\usepackage{palatino}%    Choose default roman font.  Others are times, pslatex, newcent, bookman, chancery
\usepackage{mathpazo}%  Match

\usepackage{helvet}%      Choose default sans serif

\usepackage{simplemargins}
\usepackage{slashed}
\usepackage{doi}
\usepackage{comment}

\newcommand{\MarginText}[1]{\section{\textit{#1}}}

\usepackage{hyperref} % Required for adding links	and customizing them
\hypersetup{colorlinks=true,
            linkcolor=darkblue,
            urlcolor=darkblue,
            citecolor=darkblue}
\usepackage{ifthen}
\newboolean{cms@italic}
\setboolean{cms@italic}{false}
\newboolean{cms@external}
\setboolean{cms@external}{false}
\newboolean{extended}
\setboolean{extended}{false}


\newcommand*\publistbasestyle{phys}
\usepackage[style=publist,
biblabel=brackets,
sorting=ddt,
plauthorhandling=highlight,
nameorder=given-family,
]{biblatex}
 \plauthorname[Javier]{Duarte}

\DeclareSourcemap{
 \maps[datatype=bibtex,overwrite=true]{
  \map{
    \step[fieldsource=Collaboration, final=true]
    \step[fieldset=usera, origfieldval, final=true]
  }
 }
}

\renewbibmacro*{author}{%
  \iffieldundef{usera}{%
    \printnames{author}%
  }{%
    \printfield{usera} Collaboration% \printnames{author}%
  }%
}%

\setcounter{biburllcpenalty}{7000}
\setcounter{biburlucpenalty}{8000}
% Name of your .bib file(s)
\addbibresource{bib_publications.bib}
\addbibresource{bib_reviews.bib}
\addbibresource{bib_refproceedings.bib}
\addbibresource{bib_proceedings.bib}
\addbibresource{bib_workinprogress.bib}
\addbibresource{bib_bookchapters.bib}
\addbibresource{bib_other.bib}

\newlength{\datebox}\settowidth{\datebox}{Spring 2011} % Set the width of the date box in each block

\newcommand{\NewEntry}[3]{\noindent\hangindent=2em\hangafter=0 \parbox{\datebox}{\small \textit{#1}}\hspace{1.5em} #2 #3 % Define a command for each new block - change spacing and font sizes here: #1 is the left margin, #2 is the italic date field and #3 is the position/employer/location field
\vspace{0.5em}} % Add some white space after each new entry

\newcommand{\Description[1]}{\noindent #1\vspace{-2pt}}

\usepackage[pazoGreek]{heppennames2}
\usepackage{ptdr-definitions}

\newsectionwidth{10pt} % Stops section indenting

\setallmargins{1in}

\begin{document}

%----------------------------------------------------------------------------------------
%	NAME AND CONTACT INFORMATION SECTION
%----------------------------------------------------------------------------------------

\address{{\bf Javier M. Duarte~\orcidicon{0000-0002-5076-7096}}~\inspireicon{1054778}~\webofsciicon{AAA-5414-2020}~\arxivicon{duarte_j_1} \\
  Department of Physics, 0319\\
  University of California San Diego\\
  9500 Gilman Drive\\
  La Jolla, CA 92093-0319
}

\address{\\
  office: Mayer Hall Addition 5513\\
  email: \href{mailto:jduarte@ucsd.edu}{jduarte@ucsd.edu}\\
  web: \href{https://jduarte.physics.ucsd.edu/}{https://jduarte.physics.ucsd.edu}\\
  phone: (858) 246-4980
}

\newcommand{\DOI}[1]{\href{https://doi.org/#1}{doi:#1}}
\begin{resume}

  %----------------------------------------------------------------------------------------
  %	EDUCATION
  %----------------------------------------------------------------------------------------

  \MarginText{Education}
  \noindent Ph.D. Physics, California Institute of Technology\hfill 2016
  \begin{itemize}
    \itemsep-0.3em
    \item[] Thesis: \href{https://doi.org/10.7907/Z9GX48JV}{Naturalness confronts nature: Searches for supersymmetry with the CMS detector in pp collisions at $\sqrt{s}$ = 8 and 13\TeV}
    \item[] Advisor: Maria Spiropulu
  \end{itemize}
  \vspace{-1.3em}
  \noindent M.S. Physics, California Institute of Technology\hfill 2015\\
  \noindent S.B. Physics and Mathematics, Massachusetts Institute of Technology\hfill 2010
  \begin{itemize}
    \itemsep-0.3em
    \item[] Thesis: \href{http://hdl.handle.net/1721.1/61255}{Exotic antineutrino oscillations ($\bar\nu_\Pe\to\bar\nu_{\slashed{\Pe}}$) in Double Chooz}
    \item[] Advisor: Janet Conrad
  \end{itemize}

  %----------------------------------------------------------------------------------------
  %	PROFESSIONAL EXPERIENCE
  %----------------------------------------------------------------------------------------

  \MarginText{Professional Experience}
  Assistant Professor of Physics at UC San Diego, La Jolla, CA \hfill 2019--Present\\
  \href{http://www.fnal.gov/pub/forphysicists/fellowships/leon_lederman/index.html}{Lederman Fellow} at Fermilab, Batavia, IL \hfill 2016--2019\\
  Technical Instructor in \href{http://web.mit.edu/8.13/www/index.shtml}{Junior Lab} at MIT, Cambridge, MA \hfill 2010--2011

  \MarginText{Grants and External Funding}
  \begin{itemize}
    \itemsep-0.3em
    \item Key Personnel and Institute PI for \href{https://www.nsf.gov/awardsearch/showAward?AWD_ID=2117997&HistoricalAwards=false}{NSF HDR Institute for Accelerated AI Algorithms for Data Driven Discovery (A3D3)} ({2021--2026}).
    \item Key Personnel for DOE Award for \href{https://science.osti.gov/-/media/hep/pdf/Awards/HEP_instrumentation_awards_list.pdf}{HEP Consortium for Advanced Training (HEPCAT)} ({2021--2024}).
    \item Co-PI for DOE Award for \href{https://science.osti.gov/-/media/ascr/pdf/funding/2021/Data_Reduction_Awards.pdf}{Real‐time Data Reduction Codesign at the Extreme Edge for Science} ({2021--2024}).
          % Key Personnel for NSF \href{https://www.nsf.gov/awardsearch/showAward?AWD_ID=2112606&HistoricalAwards=false}{AI Institute for Intelligent CyberInfrastructure with Computational Learning in the Environment (ICICLE)} ({2021--2026}).
    \item PI of DOE Early Career Award for \href{https://science.osti.gov/-/media/early-career/pdf/FY20_DOE_SC_Early_Career_Research_Program_Abstracts.pdf}{Real-Time Artificial Intelligence for Particle Reconstruction and Higgs Physics} ({2020--2025}).
    \item Co-PI of DOE Award for \href{https://science.osti.gov/-/media/ascr/pdf/programdocuments/docs/2020/List_of_Awards_FAIR_Data.pdf}{FAIR Framework for Physics-Inspired Artificial Intelligence in High Energy Physics} ({2020--2023}).
    \item Co-PI of NSF Award for \href{https://nsf.gov/awardsearch/showAward?AWD_ID=2005369}{Exploring Neural Network Processors for AI in Science and Engineering} ({2020--2021}).
    \item Key Personnel for \href{https://www.internet2.edu/news/detail/17957/}{Investigating Heterogeneous Computing at the Large Hadron Collider} Phase-II sub-award of Internet2 NSF Grant \href{https://www.nsf.gov/awardsearch/showAward?AWD_ID=1904444}{Exploring Clouds for Acceleration of Science (E-CAS)} ({2020-2021}).
    \item DOE QuantISED Award \href{https://pamspublic.science.energy.gov/WebPAMSExternal/Interface/Common/ViewPublicAbstract.aspx?rv=1f7d4729-6f93-40bd-a55f-c108545b1ea9&rtc=24&PRoleId=10}{Quantum Machine Learning and Quantum Computation Frameworks for HEP (QMLQCF)} ({2018--2020}).
    \item Fermilab LDRD Award for \href{https://ldrd.fnal.gov/subdir/FNAL-LDRD-2019-017-D1.pdf}{Graph Neural Networks for Accelerating Calorimetry and Event Reconstruction} ({2019--2021}).
    \item Fermilab LDRD Award for \href{https://ldrd.fnal.gov/subdir/FNAL-LDRD-2019-027-D1.pdf}{Accelerator Control with Artificial Intelligence} ({2019--2021}).
  \end{itemize}
  %----------------------------------------------------------------------------------------
  %	FELLOWSHIPS AND AWARDS
  %----------------------------------------------------------------------------------------

  \MarginText{Fellowships and Awards}
  \begin{itemize}
    \itemsep-0.3em
    \item UCSD Undergraduate Research Hub \href{https://ugresearch.ucsd.edu/mentors/outstanding-mentor-award/index.html#2021}{Outstanding Mentor Award} ({2021})
    \item \href{https://science.osti.gov/-/media/early-career/pdf/FY20_DOE_SC_Early_Career_Research_Program_Abstracts.pdf}{DOE Early Career Award} ({2020})
    \item William A. Lee Chancellor's Endowed Junior Faculty Fellow II ({2019--Present})
    \item \href{http://lpc.fnal.gov/fellows/2019/Javier_Duarte.shtml}{LHC Physics Center Distinguished Researcher} ({2019})
    \item \href{http://www.fnal.gov/pub/forphysicists/fellowships/leon_lederman/index.html}{Fermilab Lederman Fellowship} ({2016--2019})
          \ifthenelse{\boolean{extended}}{
    \item \href{http://pma.caltech.edu/research-and-academics/physics/physics-prizes-awards/france-a-cordova-graduate-student-fund-recipients}{Tombrello Scholar} sponsored by the France A. C\'{o}rdova Graduate Student Fund ({2015})
          }{}
    \item \href{https://www.nsfgrfp.org/}{NSF Graduate Research Fellowship} ({2011--2014})
          %\Description{{2012} \hspace{1.5em} Isidor I. Rabi Award at the \href{http://www.emfcsc.infn.it/issp2012/index.html}{International School of Subnuclear Physics}}\\
    \item \href{http://www.gmsp.org/}{Gates Millenium Scholar} sponsored by Hispanic Scholarship Fund ({2006--2014})
  \end{itemize}

  %----------------------------------------------------------------------------------------
  %	Research
  %----------------------------------------------------------------------------------------

  \MarginText{Selected Research Experience}

  \textbf{Higgs Boson Measurements \& Combinations}
  \begin{itemize}
    \itemsep-0.3em
    \item Author of CMS search for high-momentum (boosted) double Higgs boson production in the four bottom quark final state~\cite{CMS:2022nmn} and CMS statistical combination of Higgs boson pair searches~\cite{CMS:2022dwd} using full Run 2 data ({2020--Present}).
    \item Lead author of search for a highly boosted Higgs boson decaying to a bottom quark-antiquark pair using full Run 2 data~\cite{Sirunyan:2020hwz} and contributions to the corresponding charm quark-antiquark search~\cite{CMS-PAS-HIG-21-012} ({2017--Present}).
    \item Co-author of first search for a highly boosted Higgs boson decaying to a bottom quark-antiquark pair using 2016 data, published in \emph{Phys. Rev. Lett.}~\cite{Sirunyan:2017dgc}; Adapted analysis for interpretation for differential gluon fusion Higgs boson $\pt$ measurement~\cite{Sirunyan:2018sgc}; Combination of this result with other channels led to an observation of $\PH(\bbbar)$ decay~\cite{Sirunyan:2018kst} and other measurements~\cite{Sirunyan:2018koj} ({2016--2018}).
    \item Level-1 and high-level (software) trigger development for Higgs decaying to bottom quark-antiquark pairs produced in association with a $\PZ$ boson decaying to neutrinos~\cite{Sirunyan:2018kst} ({2016--2017}).
  \end{itemize}

  \textbf{Exotic Long-lived Particle and Dijet Searches}
  \begin{itemize}
    \itemsep-0.3em
    \item Supervision of students and postdoctoral researchers performing searches for long-lived particles~\cite{CMS:2021yhb,CMS:2021juv} ({2019--Present}).
    \item Co-convener of CMS Exotica Jets+X subgroup ({2018--2020}).
    \item Co-leader of analysis group searching for exotic, light spin-$1$ and spin-$0$ particles decaying to quarks~\cite{Sirunyan:2019vxa,Sirunyan:2019sgo,Sirunyan:2018ikr,Sirunyan:2017nvi} ({2017--2019}).
    \item Co-leader of dijet resonance search group, including data scouting, wide resonance, and $\PQb$-tagged resonance searches~\cite{Sirunyan:2019pnb,Sirunyan:2019vgj,CMS-PAS-EXO-17-026,Duarte:2018bsd,Sirunyan:2018xlo,Sirunyan:2016iap} ({2016--2018}).
  \end{itemize}

  \textbf{Novel Machine Learning Algorithms for Physics}
  \begin{itemize}
    \itemsep-0.3em
    \item Development of frameworks for sharing findable, accessible, interoperable, and reusable (FAIR) data and models in high energy physics~\cite{Huerta:2022kgj,Chen:2021euv} ({2020--Present}).
    \item Development of anomaly detection algorithms for new physics searches~\cite{Govorkova:2021utb,Jawahar:2021vyu,Aarrestad:2021oeb,Tsan:2021brw,Kasieczka:2021xcg,Wozniak:2020} ({2019--Present}).
    \item Development of graph neural networks~\cite{Duarte:2020ngm} for particle-flow reconstruction~\cite{Pata:2022wam,CMS-DP-2021-030,Pata:2021oez}, including explainable AI techniques~\cite{Mokhtar:2021bkf}, and charged particle tracking~\cite{Elabd:2021lgo,Dezoort:2021kfk,Heintz:2020soy} ({2019--Present}).
    \item Supervision of students developing generative adversarial networks and autoencoders for fast sparse data generation in high energy physics~\cite{Touranakou:2022qrp,Kansal:2021cqp,Orzari:2021suh,Kansal:2020svm} ({2019--Present}).
    \item Contributor to the Snowmass 2022 Community Planning Exercise, including white papers on machine learning for Higgs boson pair production~\cite{Apresyan:2022tqw}, graph neural networks~\cite{Thais:2022iok}, fast machine learning~\cite{Harris:2022qtm}, and data science and machine learning in physics education~\cite{Benelli:2022sqn}; Co-convenor of the CompF4 subgroup on AI Hardware~\cite{Bhimij:2022xyn}.
    \item Studies of quantum machine learning and quantum computation frameworks for high energy physics, including charged particle tracking~\cite{Zlokapa:2019tkn} ({2018--Present}).
    \item Development of deep neural networks for identifying boosted Higgs bosons decaying to $\bbbar$ and $\ccbar$ inside and outside of  CMS~\cite{Moreno:2019neq,Moreno:2019bmu,neurips2019_hbb,CMS-DP-2018-046} ({2017--Present}).
  \end{itemize}

  \textbf{Fast Machine Learning Inference for Physics}
  \begin{itemize}
    \itemsep-0.3em
    \item Fast machine learning scientific benchmarks~\cite{Duarte:2022hdp} (2022--Present).
    \item Development of Quantized ONNX (QONNX) framework for representing arbitrary-precision neural networks~\cite{Pappalardo:2022nxk} (2022--Present).
    \item Real-time AI on FPGAs for accelerator control~\cite{John:2020sak} ({2018--2021}).
    \item Development \texttt{hls4ml} for scientific low-power machine learning devices~\cite{Hussain:2022faststamp,Borras:2022opensource,Banbury:2021mlperf,DiGuglielmo:2021ide,Hawks:2021ruw,Fahim:2021cic} ({2020--Present}).
    \item Development of \href{https://github.com/fastmachinelearning/SonicCMS}{Services for Optimal Network Inference on Coprocessors (SONIC)}~\cite{Krupa:2020bwg,Rankin:2020usv,neurips2019_sonic,Duarte:2019fta} ({2018--Present}).
    \item Creation of \href{https://fastmachinelearning.org/hls4ml/}{\texttt{hls4ml}} tool for creating low-latency FPGA-based firmware implementations of machine learning algorithms~\cite{Deiana:2021niw,Elabd:2021lgo,Govorkova:2021utb,Iiyama:2020wap,Heintz:2020soy,DiGuglielmo:2020eqx,Summers:2020xiy,neurips2019_hls4ml,Duarte:2018ite} ({2017--Present}).
    \item R\&D, including firmware development and hardware demonstration, for the CMS Global Correlator Trigger for the Phase-2 upgrade of the Level-1 trigger~\cite{CERN-LHCC-2020-004} (\emph{2017--Present}).
  \end{itemize}

  \ifthenelse{\boolean{extended}}{
    \textbf{Trigger.} Tested 3D Vertically-Integrated Pattern Recognition Associate Memory (VIPRAM) chips for the Phase-2 upgrade of the CMS Level-1 track trigger ({2016--2017}).

    \textbf{SUSY Analysis.} Lead analyzer for supersymmetry (SUSY) searches using razor variables at 8 and 13\TeV~\cite{Khachatryan:2016epu,Duarte:2016wnw,Khachatryan:2015pwa,Duarte:2014soa} ({2012--2016}).

    \textbf{SUSY Phenomenology.} Development of a simplified model~\cite{Duarte:2017bbq} for bottom-squark-mediated Higgs+jet production used to interpret a CMS search for supersymmetry with $\PH\to\gamma\gamma$ decays~\cite{Sirunyan:2017eie} ({2015--2016}).

    \textbf{Timing Detector.} Measurements of the timing properties of LYSO-based calorimeters at Fermilab and CERN test beam facilities, demonstrating that a timing resolution of approximately 30\unit{ps} is achievable~\cite{Bornheim:2017gql,8069874,Anderson:2016tiu,Anderson:2016ygg,7581887,Anderson:2015gha,Anderson:2015tia,Bornheim_2015} ({2014--2016}).

    \textbf{SUSY Trigger Algorithms.} Development and implementation of high-level trigger algorithms for selecting SUSY and dark matter candidate events using razor variables for Run 2 of the LHC at 13 TeV  ({2014--2016}).

  }{}

  %----------------------------------------------------------------------------------------
  %	PUBLICATIONS
  %----------------------------------------------------------------------------------------

  \MarginText{Selected Publications, Reports, Book Chapters, and Conference Proceedings}
  Selected publications, reports, book chapters, and conference proceedings to which I made a substantial contribution are listed here.\nocite{*}

  \vspace{11pt}
  \setlength\bibitemsep{\itemsep-0.3em}
  \printbibliography[heading=none,keyword={career}]

  %----------------------------------------------------------------------------------------
  %	TEACHING
  %----------------------------------------------------------------------------------------

  \MarginText{Teaching}

  \textbf{Computational Physics}
  \begin{itemize}
    \itemsep-0.3em
    \item Co-instructor for Physics 142/242: Computational Physics II: PDE and Matrix Models ({Spring 2022}).
    \item Co-instructor for Physics 141/241: Computational Physics I: Probabilistic Models and Simulations ({Winter 2022}).
  \end{itemize}

  \textbf{Machine Learning \& Data Science}
  \begin{itemize}
    \itemsep-0.3em
    \item Creator of and instructor for Physics 139/239: Machine Learning in Physics ({Winter 2023}).
    \item \href{https://iaifi.org/phd-summer-school.html}{NSF IAIFI Ph.D. Summer School} Lecturer on ``Representations, networks, and symmetries for learning from particle physics data'' and ``Model compression and fast machine learning in particle physics'' (Summer 2022).
    \item Particle physics domain mentor for data science capstone \href{https://jmduarte.github.io/capstone-particle-physics-domain}{DSC 180AB} ({Fall 2020, Winter 2021, Fall 2021, Winter 2022}).
    \item Guest speaker for Purdue Physics 324: Big Data Science II ({Spring 2021})
    \item Guest speaker for \href{https://github.com/violatingcp/MIT_8.S50}{MIT 8.S50}: Computational Data Science in Physics ({Winter 2021}).
    \item Creator of the \href{https://github.com/FNALLPC/machine-learning-hats}{LHC Physics Center Machine Learning Tutorial} ({2017--2020}).
  \end{itemize}

  \textbf{Computer Science \& Engineering}
  \begin{itemize}
    \itemsep-0.3em
    \item Guest speaker for \href{http://kastner.ucsd.edu/ryan/cse237c/}{CSE 237C}: Validation and Testing of Embedded Systems ({Fall 2020, Fall 2021}).
  \end{itemize}

  \textbf{Statistics for Particle Physics}
  \begin{itemize}
    \itemsep-0.3em
    \item UCSD instructor of record for \href{https://indico.cern.ch/event/1066958/}{Physics 239: Statistics in Particle Physics at the LPC} [Primary instructor: Harrison Prosper, Florida State University] ({Fall 2021}).
  \end{itemize}

  \textbf{Introductory Physics}
  \begin{itemize}
    \itemsep-0.3em
    \item Lead instructor for \href{https://jduarte.physics.ucsd.edu/phys2c/index.html}{Physics 2C}: Fluids, Waves, Thermodynamics, and Optics for 300+ undergraduate students ({Winter 2020, Winter 2021, Spring 2021}).
    \item Teaching assistant in statistical and quantum mechanics at Caltech ({2011--2012}).
  \end{itemize}

  \textbf{Seminar Courses}
  \begin{itemize}
    \itemsep-0.3em
    \item Organizer of \href{https://indico.cern.ch/event/956641/}{Physics 191}: Undergraduate Seminar on Physics ({Fall 2020}) and guest speaker ({Fall 2019}).
    \item Guest speaker for Physics 261: Seminar on Physics Research at UC San Diego ({Winter 2020, Winter 2021}).
    \item Guest speaker for Thurgood Marshall College 2: Transfer Year Experience. ({Fall 2021}).
    \item Facilitator for \href{http://eddata.fnal.gov/lasso/program_search/show_workshopID_new.lasso?event_id=435}{Taking Research into Your Classroom} Workshop at Waubonsee Comunity College ({2017}).
  \end{itemize}

  % \textbf{Research for Undergraduates.} Instructor for Physics 199: Research for Undergraduates for 3 students ({Winter 2021, Spring 2021}).

  \textbf{Physics Lab}
  \begin{itemize}
    \itemsep-0.3em
    \item Technical Instructor in \href{http://web.mit.edu/8.13/www/index.shtml}{MIT Junior Lab}, teaching third-year undergraduate physics students and maintaining the experiments ({2010--2011}).
  \end{itemize}

  %----------------------------------------------------------------------------------------
  %	OUTREACH
  %----------------------------------------------------------------------------------------

  \MarginText{Outreach}
  \textbf{Community Outreach}
  \begin{itemize}
    \itemsep-0.3em
    \item Speaker at \href{https://indico.cern.ch/event/1167965/timetable/#93-meet-a-uscms-professorscien}{``Meet a US CMS Professor''} for US CMS Internship Program on Wednesday, August 10 (2022).
    \item Exhibitor for UC San Diego Physics and Duarte Lab at the \href{https://www.barriologansae.com/}{Barrio Logan Science \& Art Expo} on Saturday, April 16 and \href{https://www.southeastsdsteamexpo.com/}{Souteast Science \& Art Expo} on Saturday, August 13 ({2022}).
    \item Panelist on ``Careers in STEM Teaching and Research in Higher Education Panel'' for UC San Diego Physical Sciences Division Student Success Center on Monday, November 15 ({2021}).
    \item Invited faculty speaker at UC San Diego Physics Graduate Student Diversity Initiative Grad Recruitment Fair on Saturday, October 23 ({2021}).
    \item Speaker on ``Undergraduate and graduate research opportunities'' at UC San Diego SACNAS Chapter's Community College Workshop on Saturday, November 21 ({2020}).
    \item Invited faculty speaker for \href{http://ypp.ucsd.edu/}{Young Physicists Program} at UC San Diego ({2020}).
    \item Co-director of \href{http://saturdaymorningphysics.fnal.gov/}{Saturday Morning Physics} and lectuer on \href{http://saturdaymorningphysics.fnal.gov/fall-session-2018/}{Symmetry, Antimatter, and Supersymmetry} at Fermilab ({2018--2019}).
    \item On-site coordinator for \href{http://saturdaymorningphysics.fnal.gov/}{Saturday Morning Physics} at Fermilab ({2016--2018}).
  \end{itemize}

  \textbf{Mentoring}
  \begin{itemize}
    \itemsep-0.3em
    \item Mentor for Cal-Bridge program ({2021--2023}).
    \item Mentor for \href{http://resilientmaterials.ucsd.edu/ENLACE}{ENLACE binational summer research program} ({2021--2022}).
    \item Mentor in U.S. CMS Mentorship Program ({2020--2022}).
  \end{itemize}

  \textbf{Advocacy}
  \begin{itemize}
    \itemsep-0.3em
    \item U.S. LHC Users Association Annual Trip to U.S. Congress ({2017, 2021}).
  \end{itemize}

  \textbf{Diversity Programs}
  \begin{itemize}
    \itemsep-0.3em
    \item Creator and coordinator of A3D3 Postbaccalaureate Fellowship Program (2021--Present).
    \item Mentor in the \href{http://diversity.fnal.gov/sist/}{SIST} internship program at Fermilab ({2018})
    \item Graduate student ambassador for the \href{http://diversity.fnal.gov/fshpe/}{Fermilab SHPE chapter} ({2018}).
    \item Member of the \href{http://diversity.fnal.gov/target/}{TARGET} program committee at Fermilab ({2017--2019}).
    \item Residential Facilitator for \href{http://ome.mit.edu/programs-services/program-overview}{MIT Interphase EDGE program} ({2010}).
  \end{itemize}

  %----------------------------------------------------------------------------------------
  %	PRESENTATIONS
  %----------------------------------------------------------------------------------------

  \MarginText{Selected Conference, Workshop, and Seminar Presentations}

  \begin{itemize}
    \itemsep-0.3em
    \item \href{https://indico.cern.ch/event/1211501/contributions/5128872/}{Recent advances in machine learning for high energy physics}. Dark Interactions: New perspectives from Theory and Experiment. November 16, 2022.
    \item \href{https://indico.cern.ch/event/1219204/}{Measuring Higgs bosons using artificial intelligence}. Physics Department Colloquium. November 3, 2022. University of California San Diego, La Jolla, CA, USA.
    \item \href{https://www.sandiego.edu/events/cas/detail.php?_focus=86860}{Measuring Higgs bosons using artificial intelligence}. The Human Side of Science Lecture Series. November 1, 2022. University of San Diego, San Diego, CA, USA.
    \item \href{https://iaifi.org/summer-school-agenda}{Model compression and fast machine learning in particle physics}. IAIFI Summer School. August 2, 2022. Tufts University, MA, USA.
    \item \href{https://iaifi.org/summer-school-agenda}{Representations, networks, and symmetries for learning from particle physics data}. IAIFI Summer School. August 1, 2022. Tufts University, Medford, MA, USA.
    \item \href{https://indico.fnal.gov/event/22303/contributions/245447}{CompF3: ML for Data Analysis Summary}. Community Summer Study Snowmass 2022. July 19, 2022. University of Washington, Seattle, WA, USA.
    \item \href{https://indico.fnal.gov/event/22303/contributions/245914}{CompF4: AI-Hardware Summary}. Community Summer Study Snowmass 2022. July 19, 2022. University of Washington, Seattle, WA, USA.
    \item \href{https://indico.fnal.gov/event/22303/contributions/244766}{A3D3 Postbaccalaureate Fellowship Program}. Community Summer Study Snowmass 2022. July 18, 2022. University of Washington, Seattle, WA, USA.
    \item \href{https://indico.fnal.gov/event/55499/}{Enabling the Higgs self-coupling measurement with highly energetic Higgs pairs in CMS}. Joint Experimental-Theoretical Physics Seminar (Wine \& Cheese). June 3, 2022. Fermilab, Batavia, IL, USA.
    \item \href{https://indico.cern.ch/event/1128328/contributions/4900731/}{Accelerated Graph Neural Network Inference}. Mini-workshop on Graph Neural Networks for Tracking Colocated with Connecting the Dots 2022. June 3, 2022.
    \item \href{https://sites.google.com/g.harvard.edu/mlperf-bench-hpca22/home}{Fast Machine Learning for Science}. ML Performance: Benchmarking Deep Learning Systems (MLPerf-Bench) Tutorial at the 28th IEEE International Symposium on High-Performance Computer Architectures (HPCA 2022). April 3, 2022.
    \item \href{https://physics.columbia.edu/events/hep-seminar-dr-javier-duarte-ucsd}{AI at the Edge of Particle Physics}. HEP Seminar. November 17, 2021. Columbia University, New York, NY, USA.
    \item \href{https://www.anl.gov/event/ai-at-the-edge-of-particle-physics}{AI at the Edge of Particle Physics}. AI Distinguished Lecture Series. August 12, 2021. Argonne National Laboratory.
    \item \href{https://tdai.osu.edu/events/acm-symposium-harnessing-data-science-autonomous-computing-materials}{AI-Hardware Codesign for Real-Time Science}. Harnessing Data Science for Autonomous Computing Materials Symposium. May 27, 2021. Ohio State University.
    \item \href{https://www.calstatela.edu/dept/physics/colloquia-and-events}{The Edge of Particle Physics}. Department of Physics and Astronomy Colloquium. May 6, 2021. Cal State LA.
    \item \href{https://cse.umn.edu/physics/elementary-particle-physics-elem-part-phys-seminar}{Graph Neural Networks for High Energy Physics}. Elementary Particle Physics Seminar. April 28, 2021. University of Minnesota.
    \item \href{https://www.tinyml.org/event/research-symposium-2021/}{\texttt{hls4ml}: An open-source codesign workflow to empower scientific low-power machine learning devices}. tinyML Research Symposium. March 26, 2021.
    \item \href{https://indico.cern.ch/event/955026/}{Graph neural network tracking on FPGAs}. IRIS-HEP Topical Meeting. October 21, 2020.
    \item \href{https://www.cs.cmu.edu/calendar/fri-2020-10-16-1200/ece-graduate-seminar}{Real-time AI in particle physics}. ECE Graduate Seminar. October 16, 2020. Carnegie Mellon University, Pittsburgh, PA, USA.
    \item \href{https://indico.fnal.gov/event/23109/contributions/193292/}{Recent highlights from CMS}. 53rd Annual Fermilab Users Meeting. August 12, 2020. Fermilab, Batavia, IL, USA.
    \item \href{https://indico.fnal.gov/event/22961/}{Deep learning for Higgs and new physics at the LHC}. High Energy Physics Division Seminar. April 8, 2020. Argonne National Laboratory, Lemont, IL, USA.
    \item \href{https://ml4physicalsciences.github.io/2019/files/NeurIPS_ML4PS_2019_74.pdf}{``Low-latency machine learning inference on FPGAs.''} 2nd Machine Learning and the Physical Sciences Workshop at NeurIPS 2019.  December 14, 2019. Vancouver, Canada.
    \item \href{https://physics.drupal.ku.edu/calendar/colloquia#/?i=2}{Deep learning for Higgs and new physics at the LHC}. Department of Physics and Astronomy Colloquium. November 11, 2019. University of Kansas, Lawrence, KS, USA.
    \item \href{http://physics.bu.edu/events/show/2204}{Deep learning for Higgs and new physics at the LHC}. High Energy Experiment Seminar. October 10, 2019. Boston University, Boston, MA, USA.
    \item \href{https://escience2019.sdsc.edu/program}{Machine learning on FPGAs for low-latency and high-throughput inference}. eScience 2019. September 24-27, 2019. San Diego, CA, USA.
    \item \href{https://indico.cern.ch/event/793125/contributions/3495251/}{Deep learning on FPGAs tutorial}. 1st Real Time Analysis Workshop. July 15-26, 2019. Institute Pascal, Universit\'{e} Paris-Saclay, Saint Aubin, France.
    \item \href{https://indico.cern.ch/event/687651/contributions/3428206/}{Machine learning using CERN open data}. LHCP 2019. May 20-25, 2019. Benem\'{e}rita Universidad Aut\'{o}noma de Puebla, Puebla, Mexico.
    \item \href{https://indico.cern.ch/event/687651/contributions/3426898/}{Dark sector searches in CMS}. LHCP 2019. May 20-25, 2019. Benem\'{e}rita Universidad Aut\'{o}noma de Puebla, Puebla, Mexico.
    \item \href{https://indico.cern.ch/event/742793/contributions/3274392/}{FPGA-accelerated machine learning inference for particle physics}. Connecting the Dots 2019. April 2, 2019. Valencia, Spain.
    \item \href{https://universityofchicago.hosted.panopto.com/Panopto/Pages/Viewer.aspx?id=66ca09d7-74c1-4b12-bb57-a9fa01046cdf}{Unlocking the potential of LHC data: boosted Higgs and deep learning}. Particle Physics Special Seminar. February 20, 2019. University of Chicago, Chicago, IL, USA.
    \item \href{http://theory.fnal.gov/events/event/results-from-cms-18/}{Boosted Higgs couplings and dark mediators with deep learning in CMS}. Joint Experimental-Theoretical Physics Seminar (Wine \& Cheese). December 14, 2018. Fermilab, Batavia, IL, USA.
    \item \href{https://indico.cern.ch/event/745718/contributions/3211982/}{Heavy flavour identification for boosted resonances and large cone jets in CMS}. Machine Learning for Jet Physics (ML4Jets) 2018. November 14-16, 2018. Fermilab, Batavia, IL, USA.
    \item \href{https://www.physicsandastronomy.pitt.edu/events/hep-seminar-javier-duarte-fermilab}{Boosted Higgs, dark matter, and deep learning}. High Energy Physics Seminar. October 3, 2018. University of Pittsburgh, Pittsburgh, PA, USA.
    \item \href{https://indico.cern.ch/event/697988/contributions/3055990/}{\texttt{hls4ml}: Deploying Deep Learning on FPGAs for L1 trigger and Data Acquisition}. Topical Workshop on Electronics for Particle Physics (TWEPP) 2018. September 17-21, 2018. KU Leuven Campus Carolus, Antwerp, Belgium.
    \item \href{https://conferences.lbl.gov/event/137/session/27/contribution/354}{Searches for Dark Matter Mediators with the CMS Detector}. Conference on the Intersections of Particle and Nuclear Physics (CIPANP) 2018. May 29 - June 3, 2018. Hyatt Regency Indian Wells Conference Center, Indian Wells, CA, USA.
    \item \href{https://indico.fnal.gov/event/16908/}{Fast inference of deep neural networks in FPGAs for particle physics}. Research Techniques Seminar. April 24, 2018. Fermilab, Batavia, IL, USA.
    \item \href{https://indico.cern.ch/event/658267/contributions/2881127/}{Fast reconstruction and data scouting}. Connecting the Dots 2018. March 20-22, 2018. University of Washington, Seattle, WA, USA.
    \item \href{http://www.thphys.uni-heidelberg.de/~higgs/talks/duarte.pdf}{Boosted Higgs in CMS}. Higgs Couplings 2017. November 6-10, 2017. Heidelberg University, Heidelberg, Germany.
    \item \href{https://indico.hep.caltech.edu/indico/conferenceDisplay.py?confId=149}{Unlocking the potential of CMS data: boosted Higgs, low-mass dijet resonances, and data scouting}. High Energy Physics Seminar. October 30, 2017. Caltech, Pasadena, CA, USA.
    \item \href{https://indico.cern.ch/event/615891/contributions/2666361/}{Search for low-mass dijet resonances}. TeVPA 2017. August 7-11, 2017. Columbus, OH, USA.
    \item \href{https://indico.fnal.gov/contributionDisplay.py?sessionId=14&contribId=38&confId=11999}{Inclusive search for boosted SM Higgs bosons using H to bb decays with the CMS detector at 13 TeV}. APS DPF 2017. July 31 - August 4, 2017. Fermilab, Batavia, IL, USA.
    \item \href{https://indico.cern.ch/event/649575/}{Inclusive Higgs boson search using $\PH\to\bbbar$ decays}. Collider Cross Talk. July 20, 2017. CERN, Geneva, Switzerland.
    \item \href{https://indico.fnal.gov/contributionDisplay.py?sessionId=6&contribId=54&confId=13497}{Introduction to CMS open data for boosted object tagging with machine learning applications}. Data Science at High Energy Physics (DS@HEP) 2017. May 8-12, 2017. Fermilab, Batavia, IL, USA.
          \ifthenelse{\boolean{extended}}{
    \item \href{https://indico.cern.ch/event/540843/contributions/2464658/}{Dijet and boosted dijet searches for low-mass resonances}. New Physics Interpretations at the LHC 2. April 5-7, 2017. Argonne National Laboratory, Lemont, IL, USA.
    \item \href{https://indico.cern.ch/event/592671/contributions/2401900/}{Data scouting and boosted dijets: searches for low-mass dijet resonances}. Digging Deeper at LHC Run II. February 23-25, 2017. University of Pittsburgh, Pittsburgh, PA, USA.
    \item \href{https://indico.cern.ch/event/443176/contributions/2148316/}{Inclusive searches for SUSY using the razor variables in CMS}. SUSY 2016. July 3-8, 2016. The University of Melbourne, Melbourne, Australia.
    \item \href{https://indico.cern.ch/event/443176/contributions/2154549/}{Searches for BSM physics in dijet and multijet final states at CMS}. SUSY 2016. July 3-8, 2016. The University of Melbourne, Melbourne, Australia.
    \item \href{https://indico.cern.ch/event/331032/contributions/1720249/}{Third generation SUSY searches with the CMS detector}. SUSY 2015. August 23-29, 2015. Lake Tahoe, CA, USA.
    \item \href{https://indico.ific.uv.es/indico/contributionDisplay.py?sessionId=24&contribId=289&confId=2025}{Inclusive SUSY searches at CMS}. ICHEP 2014. July 2-9, 2014. Valencia Conference Centre, Valencia, Spain.
    \item \href{https://indico.cern.ch/event/279518/contributions/634785/}{Search for SUSY with the razor variables at CMS}. LHCP 2014. June 2-7, 2014. Columbia University, New York, NY, USA.
    \item \href{https://indico.cern.ch/event/252857/contributions/1579321/}{Searches for Supersymmetry with the CMS Detector}. 2nd High Energy Physics in the LHC Era. December 16-20, 2013. Universidad T\'{e}cnica Federico Santa Mar\'{i}a, Valpara\'{i}so, Chile.
    \item \href{https://indico.cern.ch/event/261650/contributions/586374/}{Razor Search for SUSY in CMS at $\sqrt{s}=8\TeV$}. Kinematic Variables for New Physics at the LHC. October 4-6, 2013. Caltech, Pasadena, CA, USA.
    \item \href{https://indico.cern.ch/event/218693/contributions/1520333/}{BSM Likelihoods in CMS}. Likelihoods for the LHC Searches Workshop. January 21-23, 2013. CERN, Geneva, Switzerland.
          }{}
  \end{itemize}

  %----------------------------------------------------------------------------------------
  %	COMMITTEES
  %----------------------------------------------------------------------------------------

  \MarginText{Service and Committee Work}

  \begin{itemize}
    \itemsep-0.3em
    \item Faculty advisor for UC San Diego SACNAS Chapter ({2022}).
    \item Coordinator of A3D3 NSF Institute Equity \& Career Committee, including Postbaccalaureate Fellowship Program ({2022}).
    \item US CMS Collaboration Annual Meeting Planning Committee ({2021--2022}).
    \item US CMS Collaboration Diversity, Equity, and Inclusion Committee ({2022}).
    \item Referee for \emph{Journal of High Energy Physics}, \emph{Physical Review D}, \emph{Physical Review Research}, \emph{Physics Letters B}, \emph{European Physics Journal C}, \emph{Computing and Software for Big Science}, \emph{Applied Optics}, and \emph{Nuclear Instruments and Methods in Physics Research Section A} ({2019--2022}).
    \item Reviewer for the 2022 Datasets and Benchmarks Track, and \href{https://ml4physicalsciences.github.io/2020/}{3rd} and \href{https://ml4physicalsciences.github.io/2021/}{4th} Machine Learning for the Physical Sciences Workshops at the Neural Information Processing Systems (NeurIPS) Conference ({2020--2022}).
    \item External reviewer for Swiss Data Science Center (SDSC), French National Research Agency (ANR), US Department of Energy (DOE) Early Career Research Program, and European Science Foundation (ESF) ({2019--2022}).
    \item M.S. Thesis Committee for Paul Yen Po Wang (UC San Diego) ({2021}).
    \item Faculty advisor for UC San Diego Physics Department Graduate Student Diversity Initiative ({2021}).
    \item Guest Associate Editor for \href{https://www.frontiersin.org/research-topics/19095/efficient-ai-in-particle-physics-and-astrophysics}{Efficient AI in Particle Physics and Astrophysics} Research Topic in Frontiers in Big Data and AI ({2021--2022}).
    \item UC San Diego Physics Department Equity, Diversity, and Inclusion Committee member ({2020--2022}).
    \item Scientific organizing committee for \href{https://indico.cern.ch/e/fml2020}{2nd} and \href{https://indico.cern.ch/e/fml2022}{3rd} Fast Machine Learning for Science Workshops and \href{http://www.ncsa.illinois.edu/Conferences/AcceleratedAINCSA/}{Accelerated Artificial Intelligence for Big-Data Experiments Conference}({2020--2022}).
    \item UC San Diego Physics Department Graduate Admissions Committee member ({2019--2022}).
          \ifthenelse{\boolean{extended}}{
    \item Local organizing committee for the \href{https://indico.cern.ch/e/ml4jets2018}{2018 ML4Jets} and \href{https://indico.cern.ch/e/fml}{1st Fast ML Workshop} at Fermilab LPC ({2018--2019}).
    \item \href{http://diversity.fnal.gov/target/}{Fermilab TARGET program} committee member ({2017--2019}).
    \item \href{http://lpc.fnal.gov/programs/topic/}{LPC Topic of the Week} committee chair at Fermilab ({2017--2018}).
    \item Analysis review committees for several CMS analyses: a search for new physics with dijet angular distributions~\cite{Sirunyan:2018wcm}, a search for fully hadronic $\ttbar\PH(\bbbar)$~\cite{Sirunyan:2018ygk}, a search for strongly interacting massive particles, and two searches for electroweak production of a vector-like top quark partner in the fully-hadronic final state~\cite{Sirunyan:2019xeh}.
    \item Local organizing committee for \href{http://dshep.fnal.gov}{DS@HEP 2017} workshop at Fermilab and development of public data set based on CMS open data for boosted object tagging with machine learning applications.
    \item Local organizing committee for \href{https://indico.cern.ch/event/639314/}{LHC Reinterpretation 2017} workshop at Fermilab LPC.
          }{}
  \end{itemize}

  \ifthenelse{\boolean{extended}}{
    \MarginText{Professional Memberships}

    \begin{itemize}
      \itemsep-0.3em
      \item Society for Advancement of Chicanos/Hispanics and Native Americans in Science (SACNAS) member ({2020}).
      \item National Society of Black Physicists (NSBP) member ({2020--Present}).
      \item American Physical Society (APS) member ({2020--Present}).
      \item Society of Hispanic Professional Engineers (SHPE) member ({2018--2019}).
    \end{itemize}
  }{}

  %----------------------------------------------------------------------------------------
  % STUDENTS
  %----------------------------------------------------------------------------------------

  \MarginText{Supervision and Mentorship}

  \textbf{Postdoctoral Researchers}
  \begin{itemize}
    \itemsep-0.3em
    \item Melissa Quinnan (UC San Diego). CMS level-1 trigger, Higgs boson searches ({2022--Present}).
    \item Daniel Diaz (UC San Diego). CMS level-1 trigger, long-lived particle searches ({2021--Present}).
  \end{itemize}

  \textbf{Doctoral Students}
  \begin{itemize}
    \itemsep-0.3em
    \item Zihan Zhao (UC San Diego). Self-supervised learning for jet tagging ({2022--Present}).
    \item Russell Marroquin Solares. CMS level-1 trigger long-lived particle tagger ({2022--Present}).
    \item Haoyang (Billy) Li. FAIR4HEP cookiecutter FAIR AI template ({2022--Present}).
    \item Olivia Weng (UC San Diego). Optimization of AI algorithms for FPGAs ({2021--Present}).
    \item Jieun Yoo (UIC). U.S. CMS Mentorship Program ({2021--Present}).
    \item Anthony Aportela (UC San Diego). Sloan fellow, HEPCAT fellow. Graph-based autoencoders for anomaly detection; Search for long-lived particles ({2021--Present}).
    \item Daniel Guerrero (University of Florida). U.S. CMS Mentorship Program ({2020--2021}).
    \item Farouk Mokhtar (UC San Diego). HDSI fellow, IRIS-HEP fellow. Machine learned particle-flow reconstruction; Search for boosted $\PH\to\PW\PW$ ({2020--Present}).
    \item Raghav Kansal (UC San Diego). IRIS-HEP fellow, LPC AI fellow. Graph-based generative adversarial networks for particle physics simulation; Search for boosted $\PH\PH\to\bbbar\PW\PW$ ({2019--Present}).
    \item Martin Kwok (Brown). Boosted Higgs boson search ({2018--2020}).
    \item Michael Krohn (CU Boulder). Boosted Higgs boson search, coupling measurement, and trigger development ({2017--2018}).
    \item Sean-Jiun Wang (University of Florida). Development and monitoring of triggers for the Higgs boson produced in association with a $\PZ$ boson decaying to neutrinos ({2017--2018}).
    \item Andrzej Novak (RWTH Aachen University). Development of deep neural networks for boosted Higgs identification in CMS ({2017--2019}).
    \item Jiajing Mao (Caltech). Data scouting trigger stream development ({2016--2018}).
    \item Giulia D'Imperio and Federico Preiato (Sapienza University of Rome). Dijet searches ({2016}).
  \end{itemize}

  \textbf{Masters Students}
  \begin{itemize}
    \itemsep-0.3em
    \item Nirmal Thomas (UC San Diego). Ragged batching for graph neural network inference as a service ({2022--Present}).
    \item Selwyn Reis Gomes (UC San Diego). Xilinx Alveo coprocessor support in \texttt{hls4ml} ({2022--Present}).
  \end{itemize}

  \textbf{Undergraduate Students}
  \begin{itemize}
    \itemsep-0.3em
    \item Zhaoyu Zhang (UC San Diego). Graphs GANs for particle physics ({2022--Present}).
    \item Mengke Zhang (UC San Diego). Undergraduate Research Award. Machine-learned particle-flow for the Compact Linear Collider ({2022--Present}).
    \item Anni Li (UC San Diego). IRIS-HEP fellow. Conditional generation with graph networks ({2022--Present}).
    \item Ricardo Efra\'{i}n Parra Payano (Universidad Nacional de Ingenier\'{i}a, Peru). APS National Mentoring Community ({2022}).
    \item Parvat Sapkota (University of Texas at Arlington). APS National Mentoring Community ({2021}).
    \item Saloni Agrawal (UC San Diego). EXPAND program. JetNet ({2022--Present}).
    \item Carlos Pareja (UC San Diego). EXPAND program. JetNet ({2022--Present}).
    \item Thomas Sievert (UC San Diego). FMP program. Quantum machine learning for high energy physics ({2021--Present}).
    \item Brian Sheldon (UC San Diego). FMP program. Boosted Higgs boson searches at the Future Circular Collider (hadron mode) ({2021--Present}).
    \item John Choi (UC San Diego). FMP program. Long-lived particle identification for CMS level-1 trigger ({2021--Present}).
    \item Ishaan Kavoori (UC San Diego). FMP program. FAIR4HEP cookiecutter FAIR AI template ({2021--Present}).
    \item Simon Poon (UC San Diego). FMP program. Machine-learned missing transverse momentum for CMS level-1 trigger ({2021--Present}).
    \item Sukanya Krishna (UC San Diego). IRIS-HEP fellow. Real-time anomaly detection for jets ({2021--Present}).
    \item Jason Liang (UC San Diego). tinyML with Brevitas and FINN. ({2021--Present}).
    \item Tai Nguyen (UC San Diego). Undergraduate Research Scholarship. tinyML with Brevitas and FINN. ({2021}).
    \item Han Hiller (University of Washington). UM-CERN-REU program. Machine-learned missing transverse momentum for CMS level-1 trigger ({2021}).
    \item Rohan Shenoy (UC San Diego). Undergraduate Research Award. Improved autoencoder training for HGCal ASIC for data compression ({2021--Present}).
    \item John Chen (UC San Diego). AEP program. Variable-sized-input generative graph networks ({2021}).
    \item Jevon Suharnoko (UC San Diego). Dream fellow. Transpilation of \textsc{PyTorch}-based neural networks to FPGA firmware with \texttt{hls4ml} ({2021--Present}).
    \item Rushil Roy (UC San Diego). FMP program. tinyML with \texttt{hls4ml} ({2021--Present}).
    \item Zichun Hao (UC San Diego). FMP program, Undergraduate Research Award. Lorentz-equivariant generative neural networks and $\PH\to\PW\PW$ tagging ({2021--Present}).
    \item Abdelrahman Elabd (University of Pennsylvania). IRIS-HEP fellow. Implementation of graph neural networks on FPGAs and integration into hls4ml. ({2021--2022}).
    \item Haifeng Ding (UC San Diego). FMP program. Higgs pair production sensitivity at future colliders (Snowmass study). ({2021}).
    \item Steven Tsan (UC San Diego). TRELS program. Unsupervised or semi-supervised anomaly detection algorithms for high energy physics. ({2020--Present}).
    \item Vesal Razavimaleki (UC San Diego). IRIS-HEP fellow. Implementation of graph neural networks on FPGAs. ({2019--Present}).
    \item Eric Moreno (Caltech). SURF program. Development of interaction and graph neural networks for boosted jet tagging with CMS open data. ({2018--2020}).
    \item Sydney Jenknins (University of Chicago). Compression and firmware implementation of interaction and graph neural networks for charged particle tracking at the LHC ({2018}).
    \item Eric Scotti (Brown University). Development of deep neural networks for boosted Higgs identification in CMS ({2017--2018}).
          \ifthenelse{\boolean{extended}}{
    \item Irene Li (University of Cambridge). Supervised and unsupervised machine learning methods for SUSY searches at the LHC ({2015}).
    \item Jared Filseth (Caltech). Measurement of single-lepton and razor trigger efficiencies for LHC Run 2 ({2015}).
    \item Ann Wang (Caltech). Trigger and search strategy for new Higgs-aware SUSY models at the LHC ({2013--2015}).
    \item Edward Garza (Caltech). Optimization of megajet clustering for razor SUSY searches at the LHC ({2014}).
    \item Gautam Uphadya (Caltech). New kinematic variables for gluino-mediated SUSY production ({2013--2014}).
    \item Natalie Harrison (University of Chicago). EFT Interpretation of the razor search for DM production at the LHC ({2013}).
    \item Anthony Lutz (Caltech). Generator-level validation of Monte Carlo samples for simplified SUSY models with mixed final states ({2013}).
    \item Nikita Sirohi (Caltech). Monte Carlo generation of simplified models for DM searches at LHC Run 2 ({2013}).
    \item Chun Sung Bong (Caltech). Search for MFV RPV gluinos in dileptonic final states ({2012}).
    \item Max Horton (Caltech). Frequentist and Bayesian limit-setting for razor fit-based searches ({2012}).
          }{}
  \end{itemize}

  \MarginText{Press}

  \begin{itemize}
    \itemsep-0.3em
    \item \href{https://spectrum.ieee.org/machine-learning-in-physics}{``Machine Learning Shaking Up Hard Sciences, Too''} by Dan Garisto, IEEE Spectrum, October 7, 2022
    \item \href{https://www.sciencenews.org/article/higgs-boson-particle-physics-standard-model-discovery-anniversary}{``How physicists are probing the Higgs boson 10 years after its discovery''} by Emily Conover, Science News Magazine, June 29, 2022
    \item \href{https://news.fnal.gov/2022/06/probing-higgs-self-coupling-with-boosted-higgs-pairs/}{``Probing Higgs self-coupling with boosted Higgs pairs''} by Artur Apresyan and Si Xie, Fermilab News, June 9, 2022
    \item \href{https://mlcommons.org/en/news/mlperf-inference-1q2022/}{``MLPerf Results Show Advances in Machine Learning Inference Performance and Efficiency''}, MLCommons, June 4, 2022
    \item \href{https://insidehpc.com/2022/04/mlperf-results-highlight-advances-in-machine-learning-inference-performance-and-efficiency/}{``MLPerf Results Highlight Advances in Machine Learning Inference Performance and Efficiency''}, Inside HPC, April 6, 2022
    \item \href{https://www.symmetrymagazine.org/article/double-trouble-higgs}{``Double trouble Higgs''} by Sarah Charley, Symmetry Magazine, April 26, 2022
    \item \href{https://cerncourier.com/a/graph-neural-networks-boost-di-higgs-search/}{``Graph neural networks boost di-Higgs search''}, CERN Courier, March 11, 2022
    \item \href{https://ucsdnews.ucsd.edu/pressrelease/physicists-apply-fairness-to-data-studies}{``Physicists Apply FAIRness to Data Studies''} by Kimberly Mann Bruch, UC San Diego News, February 15, 2022
    \item \href{https://www.intel.com/content/www/us/en/customer-spotlight/stories/san-diego-supercomputer-customer-story.html}{``SDSC Builds AI-Focused `Voyager' Supercomputer''}, Intel, November 9, 2021
    \item \href{https://cerncourier.com/a/muon-detector-probes-long-lived-particles/}{``Muon detector probes long-lived particles''}, CERN Courier, November 5, 2021
    \item \href{https://cerncourier.com/a/hunting-anomalies-with-an-ai-trigger/}{``Hunting anomalies with an AI trigger''}, CERN Courier, August 31, 2021
    \item \href{https://cms.cern/news/new-window-shadow-world-exotic-particle-decays-muon-detectors}{``A new window into the shadow world: Exotic particle decays in the muon detectors''}, CMS Physics Briefing, August 30, 2021
    \item \href{https://ucsdnews.ucsd.edu/pressrelease/physics-computation-experts-help-earn-15m-to-advance-ai-data-analysis}{``Physics, Computation Experts Help Earn \$15M to Advance AI, Data Analysis''} by Cynthia Dillon, UC San Diego News, September 28, 2021
    \item \href{https://ucsdnews.ucsd.edu/pressrelease/san-diego-supercomputer-center-teams-up-with-habana-to-power-voyager}{``San Diego Supercomputer Center Teams Up with Habana to Power Voyager''} by Cynthia Dillon, UC San Diego News, April 9, 2021
    \item \href{https://cms.cern/news/live-long-and-prosper-searching-long-lived-relatives-higgs-boson}{``Live long and propser: Searching for the long-lived relatives of the Higgs boson''}, CMS Physics Briefing, August 16, 2021
    \item \href{https://cerncourier.com/a/long-lived-particles-gather-interest/}{``Long-lived particles gather interest''} by James Beacham and Albert De Roeck, CERN Courier, July 21, 2021
    \item \href{https://ucsdnews.ucsd.edu/pressrelease/national-science-foundation-awards-sdsc-5-million-to-develop-innovative-ai-resource}{``National Science Foundation Awards SDSC \$5 Million to Develop Innovative AI Resource''} by Jan Zverina, UC San Diego News, July 1, 2020
    \item \href{https://cms.cern/news/boosting-unknown-highest-energy-higgs-bosons}{``Boosting into the unknown: The highest energy Higgs bosons''}, CMS Physics Briefing, May 20, 2020
    \item \href{https://physicalsciences.ucsd.edu/media-events/articles/2020/uc-san-diego-physicist-making-a-mark.html}{``UC San Diego Physicist Making a Mark''} by Cynthia Dillon, UC San Diego Division of Physical Sciences News, July 1, 2020
    \item \href{https://news.fnal.gov/2019/01/fermilab-scientists-help-push-ai-to-unprecedented-speeds/}{``Fermilab scientists help push AI to unprecedented speeds''} by Javier Duarte, Sergo Jindariani, Ben Kreis and Nhan Tran, Fermilab News, January 29, 2019
  \end{itemize}

  \MarginText{Citation Information}

  The five articles with the highest citation counts according to Google Scholar are relisted here.
  My h-index according to Web of Science is 101.
  % \newrefsection[bib_publications,sorting=none]
  \nocite{Sirunyan:2018koj}
  \nocite{Sirunyan:2016iap}
  \nocite{Duarte:2018ite}
  \nocite{Sirunyan:2018xlo}
  \nocite{Sirunyan:2019vgj}

  \vspace{11pt}
  \setlength\bibitemsep{\itemsep-0.3em}
  \printbibliography[heading=none,keyword={citation}]

  \hfill Last updated: \today

  \clearpage

  \ifthenelse{\boolean{extended}}{
    \nocite{*}
    \MarginText{All Other Publications}

    All other publications are listed here.

    \vspace{11pt}
    \printbibliography[heading=none,notkeyword={career}]
  }{}

\end{resume}

\end{document}
